% ===========================================================================
%  Chapitre 3 — Étude et représentation graphique des fonctions
%  PE 2026, composante ANALYSE
% ===========================================================================
\bmtchapitre{Étude et représentation graphique des fonctions}
  {Étudier le comportement asymptotique d'une courbe, repérer points angulés et
   points d'inflexion, situer une courbe par rapport à ses asymptotes, et
   résoudre graphiquement équations et inéquations.}
  {Analyse \textbullet\ environ 8 heures}

Vous savez maintenant calculer une limite et une dérivée. Ce chapitre ne
présente presque aucune notion nouvelle : il apprend à \emph{les faire tenir
ensemble} sur une feuille, jusqu'à produire une courbe juste.

C'est un chapitre de méthode, et c'est aussi celui qui rapporte le plus de
points au baccalauréat. Le tracé final vaut souvent un point et demi à lui
seul, et il ne s'improvise pas : chaque trait doit avoir été justifié avant
d'être posé.

% ---------------------------------------------------------------------------
\begin{revision}
Rien de neuf ici : on vérifie que les deux premiers chapitres sont en place.

\begin{enumerate}[label=\textbf{\arabic*.}, leftmargin=*, itemsep=0.35em]
  \item Déterminer $\limpinf \dfrac{3x^{2}-x+1}{x^{2}+2}$.
  \item Que signifie « $f$ est continue en $a$ » ?
  \item Comment reconnaît-on un point angulé sur une courbe ?
  \item Quel lien y a-t-il entre le signe de $f''$ et la forme de la courbe ?
  \item Donner l'ensemble de définition de $x \mapsto \dfrac{x^{2}+1}{x-3}$.
  \item Effectuer la division euclidienne de $x^{2}+1$ par $x-3$.
\end{enumerate}
\end{revision}

\begin{automatismes}
Sans calculatrice, sans brouillon.

\begin{enumerate}[label=\textbf{\alph*)}, leftmargin=*, itemsep=0.25em]
  \item $\limpinf \dfrac{2x+1}{x-4}$
  \item $\displaystyle\lim_{x \to 4^{+}} \dfrac{2x+1}{x-4}$
  \item $\limpinf \dfrac{x^{2}}{x+1}$
  \item Dérivée de $x \mapsto \dfrac{x^{2}}{x+1}$
  \item Signe de $x^{2}-4$
  \item $\limminf (x^{3}-x)$
  \item Résoudre $\dfrac{1}{x-2} > 0$
  \item Le quotient $\dfrac{x^{2}}{x+1}$ vaut-il $x-1+\dfrac{1}{x+1}$ ?
\end{enumerate}
\end{automatismes}

\begin{corrige}[automatismes]
a) $2$ \quad b) $+\infty$ \quad c) $+\infty$ \quad
d) $\dfrac{x^{2}+2x}{(x+1)^{2}}$ \quad e) négatif sur $\intervalleo{-2}{2}$,
positif ailleurs \quad f) $-\infty$ \quad g) $x > 2$ \quad
h) oui, c'est la division euclidienne de $x^{2}$ par $x+1$.
\end{corrige}

% ===========================================================================
\section{Les branches infinies}

\subsection{Asymptotes verticales et horizontales}

\begin{definition}[asymptote verticale]
Si $\displaystyle\lim_{x \to a} f(x) = \pm\infty$ — la limite pouvant n'être
qu'à gauche ou qu'à droite — alors la droite d'équation $x = a$ est
\textbf{asymptote verticale} à la courbe de $f$.
\end{definition}

\begin{definition}[asymptote horizontale]
Si $\displaystyle\lim_{x \to +\infty} f(x) = \ell$ avec $\ell$ réel, alors la
droite d'équation $y = \ell$ est \textbf{asymptote horizontale} à la courbe en
$+\infty$. Même définition en $-\infty$.
\end{definition}

\subsection{Asymptotes obliques}

\begin{definition}[asymptote oblique]
La droite d'équation $y = ax+b$, avec $a \neq 0$, est \textbf{asymptote
oblique} à la courbe de $f$ en $+\infty$ lorsque
\[
  \limpinf \bigl[f(x) - (ax+b)\bigr] = 0 .
\]
\end{definition}

\begin{methode}{trouver une asymptote oblique}
\emph{Si $f$ est une fraction rationnelle}, le plus rapide est la division
euclidienne du numérateur par le dénominateur. Écrire
$f(x) = ax+b+\dfrac{r(x)}{q(x)}$ avec $\deg r < \deg q$ : la partie entière
$ax+b$ est l'asymptote, et le reste tend vers $0$.

\smallskip
\emph{Sinon}, on cherche $a$ puis $b$ séparément :
\[
  a = \limpinf \frac{f(x)}{x}, \qquad b = \limpinf \bigl[f(x)-ax\bigr].
\]
\end{methode}

\begin{propriete}[position relative de la courbe et de l'asymptote]
Le signe de la différence $d(x) = f(x)-(ax+b)$ donne la position :
\begin{itemize}[leftmargin=1.4em, itemsep=0.15em]
  \item $d(x) > 0$ : la courbe est \emph{au-dessus} de l'asymptote ;
  \item $d(x) < 0$ : la courbe est \emph{au-dessous} ;
  \item $d(x) = 0$ : la courbe \emph{coupe} l'asymptote.
\end{itemize}
\end{propriete}

\begin{visualisation}[la courbe rejoint sa droite]
\begin{center}
\begin{tikzpicture}
\begin{axis}[
  width = 11cm, height = 6.2cm,
  axis lines = middle, xlabel = {$x$},
  xmin = -5.6, xmax = 5.6, ymin = -6.5, ymax = 6.5,
  xtick = \empty, ytick = \empty,
  axis line style = {bmtGris}, clip = true,
]
  \addplot[bmtIndigo, line width=1.3pt, domain=0.18:5.4, samples=200] {x + 1/x};
  \addplot[bmtIndigo, line width=1.3pt, domain=-5.4:-0.18, samples=200] {x + 1/x};
  \addplot[bmtLaterite, dashed, line width=0.9pt, domain=-5.4:5.4, samples=2] {x};
  \node[bmtLaterite, font=\footnotesize, anchor=west] at (axis cs:4.0,3.1) {$y=x$};
  \node[bmtGris, font=\footnotesize, anchor=west] at (axis cs:0.25,-5.4) {$x=0$};
\end{axis}
\end{tikzpicture}
\end{center}

Regardez la branche de droite : plus $x$ grandit, plus la courbe se colle à la
droite rouge sans jamais la toucher, en restant au-dessus. La branche de
gauche fait la même chose, mais par en dessous.
\end{visualisation}

\begin{erreurclassique}
Trouver $a = \limpinf \dfrac{f(x)}{x}$ fini ne suffit pas à conclure à une
asymptote. Il faut ensuite vérifier que $f(x)-ax$ admet une limite
\emph{finie}.
\end{erreurclassique}

\begin{entrainement}[branches infinies]
\exo[\niveau{1}] Déterminer les asymptotes de la courbe de
$f(x) = \dfrac{3x-1}{x+2}$.

\exo[\niveau{2}] Soit $f(x) = \dfrac{x^{2}-x+1}{x-1}$. Montrer par division
euclidienne que la droite $y = x$ est asymptote à la courbe, et étudier la
position relative.

\exo[\niveau{2}] Étudier les branches infinies de $f(x) = \dfrac{x^{3}}{x^{2}-1}$.
\end{entrainement}

\begin{corrige}
\textbf{1.} $x = -2$ est asymptote verticale ; $\lim_{x \to \pm\infty} f(x) = 3$ donc $y = 3$ est
asymptote horizontale en $\pm\infty$.

\textbf{2.} $\dfrac{x^{2}-x+1}{x-1} = x + \dfrac{1}{x-1}$. Le reste tend vers
$0$ : $y = x$ est asymptote. Son signe est celui de $x-1$ : la courbe est
au-dessus pour $x>1$, au-dessous pour $x<1$.

\textbf{3.} $\dfrac{x^{3}}{x^{2}-1} = x + \dfrac{x}{x^{2}-1}$, donc $y=x$ est
asymptote oblique. De plus $x = 1$ et $x = -1$ sont asymptotes verticales.
\end{corrige}

% ===========================================================================
\section{Ce que la courbe montre en un point}

\begin{definition}[trois points remarquables]
Soit $a$ un réel et $(\Cf)$ la courbe de $f$.
\begin{itemize}[leftmargin=1.4em, itemsep=0.25em]
  \item \textbf{Point de discontinuité} : $f$ n'est pas continue en $a$. La
        courbe est interrompue.
  \item \textbf{Point angulé} : $f$ est continue en $a$, dérivable à gauche et
        à droite, mais $f'_{g}(a) \neq f'_{d}(a)$.
  \item \textbf{Point d'inflexion} : $f''$ s'annule en $a$ \emph{en changeant
        de signe}.
\end{itemize}
\end{definition}

% ===========================================================================
\section{Le plan d'étude, une fois pour toutes}

\begin{methode}{étudier complètement une fonction}
\begin{enumerate}[label=\textbf{\arabic*.}, leftmargin=*, itemsep=0.28em]
  \item \textbf{Ensemble de définition.} Écarter les valeurs interdites.
  \item \textbf{Parité.} Si $f$ est paire ou impaire, on réduit l'étude.
  \item \textbf{Limites aux bornes.}
  \item \textbf{Branches infinies.} Asymptotes, position relative.
  \item \textbf{Dérivée et signe.} Calculer $f'$, la factoriser.
  \item \textbf{Tableau de variation.}
  \item \textbf{Points particuliers.}
  \item \textbf{Tracé.} Placer d'abord les asymptotes, puis les points connus,
        puis la courbe.
\end{enumerate}
\end{methode}

\begin{exemple}[une étude menée jusqu'au bout]
Étudions $f(x) = \dfrac{x^{2}-2x+2}{x-1}$.

\emph{Définition.} $\Df = \R \setminus \{1\}$.

\emph{Forme réduite.} $f(x) = x-1+\dfrac{1}{x-1}$.

\emph{Limites.} En $\pm\infty$, $f(x) \to \pm\infty$. En $1^{-}$ :
$f(x) \to -\infty$ ; en $1^{+}$, $f(x) \to +\infty$.

\emph{Branches infinies.} $x=1$ est asymptote verticale. Comme
$f(x)-(x-1) = \dfrac{1}{x-1} \to 0$, la droite $y = x-1$ est asymptote
oblique. La courbe est au-dessus pour $x>1$, au-dessous pour $x<1$.

\emph{Dérivée.} $f'(x) = 1 - \dfrac{1}{(x-1)^{2}} = \dfrac{x(x-2)}{(x-1)^{2}}$.

\emph{Variations.} $f$ croît sur $\left]-\infty ; 0\right]$, décroît sur
$\intervallefo{0}{1}$ puis sur $\intervalleof{1}{2}$, croît sur
$\left[2 ; +\infty\right[$. Les extremums locaux valent $f(0) = -2$ et
$f(2) = 2$.
\end{exemple}

\begin{visualisation}[l'étude précédente, achevée]
\begin{center}
\begin{tikzpicture}
\begin{axis}[
  width = 11cm, height = 6.4cm,
  axis lines = middle, xlabel = {$x$},
  xmin = -3.6, xmax = 5.6, ymin = -6.5, ymax = 6.5,
  xtick = {0,2}, ytick = {-2,2},
  axis line style = {bmtGris},
  tick label style = {font=\footnotesize, color=bmtGris}, clip = true,
]
  \addplot[bmtIndigo, line width=1.3pt, domain=1.12:5.4, samples=220] {x - 1 + 1/(x-1)};
  \addplot[bmtIndigo, line width=1.3pt, domain=-3.4:0.88, samples=220] {x - 1 + 1/(x-1)};
  \addplot[bmtLaterite, dashed, line width=0.9pt, domain=-3.4:5.4, samples=2] {x-1};
  \addplot[bmtVetiver, dashed, line width=0.9pt] coordinates {(1,-6.5) (1,6.5)};
  \filldraw[bmtIndigo] (axis cs:0,-2) circle (1.8pt);
  \filldraw[bmtIndigo] (axis cs:2,2) circle (1.8pt);
  \node[bmtLaterite, font=\footnotesize, anchor=west] at (axis cs:4.2,2.6) {$y=x-1$};
  \node[bmtVetiver, font=\footnotesize, anchor=south west] at (axis cs:1.1,-6.2) {$x=1$};
\end{axis}
\end{tikzpicture}
\end{center}

Les deux droites en pointillé ont été tracées \emph{avant} la courbe, et c'est
ce qui rend le tracé possible.
\end{visualisation}

\begin{entrainement}[études complètes]
\exo[\niveau{2}] Étudier et représenter $f(x) = \dfrac{2x^{2}-3x}{x-1}$.

\exo[\niveau{2}] Étudier et représenter $f(x) = \dfrac{x^{2}+1}{x}$, après
avoir montré que $f$ est impaire.
\end{entrainement}

\begin{corrige}
\textbf{1.} $f(x) = 2x-1-\dfrac{1}{x-1}$ ; asymptotes $x=1$ et $y=2x-1$ ;
$f'(x) = \dfrac{2x^{2}-4x+3}{(x-1)^{2}}$, de discriminant négatif donc $f'>0$ :
$f$ est croissante sur chacun des deux intervalles.

\textbf{2.} $f(-x) = -f(x)$ : impaire, courbe symétrique par rapport à
l'origine. $f(x) = x+\dfrac{1}{x}$ ; asymptotes $x=0$ et $y=x$ ;
$f'(x) = 1-\dfrac{1}{x^{2}}$, nulle en $\pm1$ : minimum $2$ en $1$, maximum
$-2$ en $-1$.
\end{corrige}

% ===========================================================================
\section{Lire et résoudre sur un graphique}

\begin{propriete}[les quatre lectures]
Soit $(\Cf)$ la courbe de $f$ et $(\mathscr{C}_{g})$ celle de $g$.
\begin{itemize}[leftmargin=1.4em, itemsep=0.22em]
  \item $f(x) = g(x)$ : abscisses des points d'intersection des deux courbes.
  \item $f(x) \leq g(x)$ : abscisses où $(\Cf)$ est au-dessous de $(\mathscr{C}_{g})$.
  \item $f(x) = m$ : abscisses des points d'intersection de $(\Cf)$ avec la
        droite horizontale d'ordonnée $m$.
  \item $f(x) \leq m$ : abscisses où $(\Cf)$ est au-dessous de cette
        horizontale.
\end{itemize}
\end{propriete}

\begin{entrainement}[lecture graphique]
\exo[\niveau{1}] La courbe de $f$ passe par $(0;2)$, admet un maximum $3$ en
$x=1$ et tend vers $-\infty$ aux deux infinis. Combien l'équation
$f(x) = 3$ a-t-elle de solutions ? Et $f(x) = 4$ ?

\exo[\niveau{2}] Deux courbes se coupent en trois points d'abscisses $-2$, $0$
et $3$, et $(\Cf)$ est au-dessus de $(\mathscr{C}_{g})$ sur
$\intervalleo{-2}{0}$ et sur $\left]3 ; +\infty\right[$. Résoudre
graphiquement $f(x) \leq g(x)$.
\end{entrainement}

\begin{corrige}
\textbf{1.} $f(x)=3$ : une solution. $f(x)=4$ : aucune solution.

\textbf{2.} L'ensemble des solutions est
$\left]-\infty ; -2\right] \cup \intervalle{0}{3}$.
\end{corrige}

% ===========================================================================
% ===========================================================================
\begin{annales}
Les études qui suivent sont reprises de sujets réellement posés ; leur
provenance est indiquée à chaque fois. Toutes portent sur des fonctions
polynômes ou rationnelles, et se mènent entièrement avec les outils des trois
premiers chapitres.

\exoann{Baccalauréat probatoire, Cameroun, série D, session 2022 — exercice 4}
Soit $f$ définie sur $\intervallefo{0}{+\infty}$ par
$f(x) = \dfrac{3x}{3+4x}$.
\begin{enumerate}[label=\textbf{\alph*)}, leftmargin=*, itemsep=0.15em]
  \item Dresser le tableau de variation de $f$.
  \item Montrer que la droite d'équation $y = \frac34$ est asymptote à la
        courbe de $f$ et préciser la position de la courbe par rapport à
        cette droite.
  \item Construire la courbe de $f$ dans un repère orthonormé.
\end{enumerate}

\exoann{Baccalauréat probatoire, Cameroun, série IH, session 2023 — problème}
Soit $f$ la fonction définie sur $\intervalle{-4}{1}$ par
$f(x) = 2x^{2}+5x-3$.
\begin{enumerate}[label=\textbf{\alph*)}, leftmargin=*, itemsep=0.15em]
  \item Dresser le tableau de variation de $f$ sur $\intervalle{-4}{1}$.
  \item Déterminer les points d'intersection de la courbe de $f$ avec les axes
        du repère.
  \item Construire cette courbe.
\end{enumerate}

\exoann{Baccalauréat probatoire, Cameroun, séries D et TI, session 2023 — exercice 4}
Soit $f$ définie par $f(x) = x+\dfrac{1}{x+1}$. Dresser le tableau de variation
de $f$, préciser ses deux asymptotes, puis construire sa courbe.

\exoann{Baccalauréat blanc, lycée Philibert Tsiranana, Terminale S, 2022-2023}
Soit $f$ définie sur $\R\setminus\{2\}$ par
$f(x) = \dfrac{x^{2}-3x+4}{x-2}$.
\begin{enumerate}[label=\textbf{\alph*)}, leftmargin=*, itemsep=0.15em]
  \item Déterminer trois réels $a$, $b$ et $c$ tels que
        $f(x) = ax+b+\dfrac{c}{x-2}$.
  \item En déduire que la droite d'équation $y = x-1$ est asymptote oblique à
        la courbe de $f$, et étudier la position relative de la courbe et de
        cette droite.
  \item Dresser le tableau de variation de $f$, puis construire sa courbe et
        ses asymptotes.
\end{enumerate}
\end{annales}

\begin{corrige}[annales]
\textbf{1. a)} $f'(x) = \dfrac{9}{(3+4x)^{2}} > 0$ : $f$ croît strictement de
$f(0) = 0$ vers $\frac34$ sur $\intervallefo{0}{+\infty}$.
\textbf{b)} $f(x)-\dfrac34 = \dfrac{-9}{4(3+4x)} < 0$ pour $x \geq 0$ : la
courbe reste sous son asymptote, et l'écart tend vers $0$.
\textbf{c)} La courbe part de l'origine, monte sans jamais atteindre le
palier $y = \frac34$.

\textbf{2. a)} $f'(x) = 4x+5$ s'annule en $-\frac54$. Sur
$\intervalle{-4}{-\frac54}$ la fonction décroît de $f(-4) = 9$ à
$f\left(-\frac54\right) = -\frac{49}{8}$ ; sur
$\intervalle{-\frac54}{1}$ elle croît jusqu'à $f(1) = 4$.
\textbf{b)} $f(0) = -3$ donne le point $(0\,;-3)$. L'équation
$2x^{2}+5x-3 = 0$ a pour racines $\frac12$ et $-3$ : la courbe coupe l'axe
des abscisses en $\left(\frac12\,;0\right)$ et $(-3\,;0)$.
\textbf{c)} Arc de parabole de sommet
$S\left(-\frac54\,;-\frac{49}{8}\right)$, tracé entre $x = -4$ et $x = 1$.

\textbf{3.} $f'(x) = \dfrac{x(x+2)}{(x+1)^{2}}$. La fonction croît sur
$\intervalleof{-\infty}{-2}$ jusqu'au maximum local $f(-2) = -3$, décroît sur
$\intervallefo{-2}{-1}$ puis sur $\intervalleof{-1}{0}$ jusqu'au minimum local
$f(0) = 1$, et croît ensuite. La droite $x = -1$ est asymptote verticale, la
droite $y = x$ asymptote oblique.

\textbf{4. a)} La division euclidienne donne
$x^{2}-3x+4 = (x-2)(x-1)+2$, d'où
$f(x) = x-1+\dfrac{2}{x-2}$ : $a = 1$, $b = -1$, $c = 2$.
\textbf{b)} $f(x)-(x-1) = \dfrac{2}{x-2} \to 0$ : la courbe est au-dessus de
l'asymptote à droite de $2$, en dessous à gauche.
\textbf{c)} $f'(x) = 1-\dfrac{2}{(x-2)^{2}}$ s'annule en $2\pm\sqrt2$. La
fonction croît jusqu'à $f\left(2-\sqrt2\right) \approx -1{,}83$, décroît puis
croît jusqu'à $f\left(2+\sqrt2\right) \approx 3{,}83$, puis croît ensuite.
\end{corrige}

\begin{jecherche}[l'atelier de briques d'Antsirabe]
Un atelier fabrique des briques. On note $q$ le nombre de milliers de briques
produites par semaine, $q \in \intervallefo{0}{+\infty}$, et l'on modélise le
bénéfice hebdomadaire, en centaines de milliers d'ariary, par
\[
  B(q) = \frac{20q}{q+2} - 4 .
\]

\begin{enumerate}[label=\textbf{\arabic*.}, leftmargin=*, itemsep=0.3em]
  \item Étudier les variations de $B$ sur $\intervallefo{0}{+\infty}$.
  \item Déterminer $\limpinf B(q)$ et en déduire que la courbe de $B$ admet
        une asymptote horizontale, dont on donnera l'équation.
  \item Interpréter économiquement cette asymptote : le bénéfice
        hebdomadaire peut-il dépasser $16$ (centaines de milliers d'ariary) ?
  \item Déterminer le seuil de rentabilité, c'est-à-dire la production $q_{0}$
        à partir de laquelle $B(q) \geq 0$.
  \item Construire la courbe de $B$ pour $q \in \intervalle{0}{20}$.
\end{enumerate}
\end{jecherche}

\begin{corrige}[l'atelier de briques d'Antsirabe]
\textbf{1.} $B'(q) = \dfrac{20(q+2)-20q}{(q+2)^{2}} = \dfrac{40}{(q+2)^{2}} > 0$ :
$B$ est strictement croissante sur $\intervallefo{0}{+\infty}$.

\textbf{2.} $\limpinf B(q) = 20-4 = 16$ : la droite $y=16$ est asymptote
horizontale.

\textbf{3.} Le bénéfice hebdomadaire augmente avec la production, mais reste
toujours strictement inférieur à $16$ (centaines de milliers d'ariary), quel
que soit le volume fabriqué : c'est un plafond que l'atelier approche sans
jamais l'atteindre.

\textbf{4.} $B(q) \geq 0 \iff 20q \geq 4(q+2) \iff 16q \geq 8 \iff q \geq
\dfrac12$ : le seuil de rentabilité est $q_{0} = 0{,}5$ millier de briques,
soit $500$ briques par semaine.

\textbf{5.} La courbe part de $B(0) = -4$, croît en se rapprochant sans
cesse de l'asymptote $y=16$, avec $B(20) = \dfrac{400}{22}-4 \approx 14{,}2$.
\end{corrige}

% ===========================================================================
\begin{jemeteste}
Répondre par vrai ou faux, en justifiant en une ligne.

\begin{enumerate}[label=\textbf{\arabic*.}, leftmargin=*, itemsep=0.28em]
  \item Une courbe ne peut pas couper son asymptote oblique.
  \item Si $\limpinf \dfrac{f(x)}{x} = 2$, la droite $y = 2x$ est asymptote à
        la courbe.
  \item Une fonction peut admettre deux asymptotes horizontales.
  \item En un point angulé, la fonction est continue.
  \item Le tracé d'une courbe doit commencer par le placement des asymptotes.
\end{enumerate}
\end{jemeteste}

\begin{corrige}[je me teste]
\textbf{1.} Faux — elle peut la couper, même plusieurs fois.
\textbf{2.} Faux — il faut encore que $f(x)-2x$ ait une limite finie.
\textbf{3.} Vrai — une en $+\infty$, une autre en $-\infty$.
\textbf{4.} Vrai — c'est dans la définition.
\textbf{5.} Vrai — c'est ce qui rend le reste du tracé fiable.
\end{corrige}

% ===========================================================================
\begin{commeaubac}
\bmtsource{Exercice au format de l'épreuve — série OSE}
Soit $f$ la fonction définie sur $\R \setminus \{-1\}$ par
$f(x) = \dfrac{2x^{2}+3x-1}{x+1}$, de courbe $(\Cf)$.

\begin{enumerate}[label=\textbf{\arabic*.}, leftmargin=*, itemsep=0.25em]
  \item Déterminer trois réels $a$, $b$ et $c$ tels que
        $f(x) = ax+b+\dfrac{c}{x+1}$. \hfill\textup{(0,75 pt)}
  \item En déduire les asymptotes de $(\Cf)$ et la position de la courbe par
        rapport à l'asymptote oblique. \hfill\textup{(1,25 pt)}
  \item Calculer $f'(x)$, étudier son signe et dresser le tableau de variation
        de $f$. \hfill\textup{(1,5 pt)}
  \item Construire $(\Cf)$ et ses asymptotes. \hfill\textup{(1,5 pt)}
  \item Discuter graphiquement, selon le réel $m$, le nombre de solutions de
        l'équation $f(x) = m$. \hfill\textup{(1 pt)}
\end{enumerate}
\end{commeaubac}

\begin{corrige}[exercice au format de l'épreuve]
\textbf{1.} La division euclidienne donne $f(x) = 2x+1-\dfrac{2}{x+1}$, donc
$a=2$, $b=1$, $c=-2$.

\textbf{2.} $x = -1$ est asymptote verticale ; $y = 2x+1$ est asymptote
oblique. Le signe de $f(x)-(2x+1) = \dfrac{-2}{x+1}$ est l'opposé de celui de
$x+1$ : la courbe est au-dessous de l'asymptote pour $x > -1$, au-dessus pour
$x < -1$.

\textbf{3.} $f'(x) = 2 + \dfrac{2}{(x+1)^{2}} > 0$ : $f$ est strictement
croissante sur chacun des deux intervalles.

\textbf{4.} La courbe possède deux branches de part et d'autre de
l'asymptote verticale.

\textbf{5.} La courbe possède deux branches, chacune parcourant $\R$ tout
entier de façon strictement croissante. Pour tout réel $m$, l'équation
$f(x)=m$ admet donc exactement deux solutions, une sur chaque branche.
\end{corrige}
